\scanpage{173}
\begin{center}- 42 -\end{center}\hfill IX

est donnée par KLEIMAN dans [14 2.8]).

a) On utilise la formule (4.2.10), qui devient ici
\[
 (4.3.0)\qquad\operatorname{Tr}g_{U_\ell}
 =\operatorname{Tr}g_{B_{\unsure{o|0}}^*,\ell}+\operatorname{Tr}g_{T_{\unsure{o|0}}^*,\ell}
 +\operatorname{Tr}g^{-1}_{T_{\unsure{o|0}}^*,\ell}\qquad,
\]
où maintenant $g_{B_{\unsure{o|0}}^*}$ et $g_{T_{\unsure{o|0}}^*}$ sont des k'-automorphismes de $B_{\unsure{o|0}}^*$ resp. de $T_{\unsure{o|0}}^*$.
Il suffit donc de montrer que les deux traces correspondantes sont
des entiers rationnels indépendants de $\ell$ ce qui est un énoncé ``géométrique'' général sur les groupes algébriques. Pour $g_{B_{\unsure{o|0}}^*}$ c'est un résultat
bien connu de WEIL sur la théorie des variétés abéliennes [5\quad][15\quad] ;
pour $g_{T_{\unsure{o|0}}^*}$ c'est pratiquement trivial, en explicitant la représentation
de $\Gamma$ sur $T=T_{\unsure{o|0}}^*$ à l'aide de la représentation correspondante sur le groupe
dual $M'$, isomorphe à $\mathbb Z^\mu$ : la trace cherchée est contragrédiente de celle
de cette dernière représentation, grâce à l'isomorphisme canonique
\[
 T_\ell(T)\simeq\check M_\ell(1)\overset{\mathrm{dfn}}{=}
 \check M'\otimes_{\mathbb Z}\mathbb Z_\ell(1)
\]

b) Si l'image de $g$ dans $\pi_{\unsure{o|0}}$ est $f_q^n$, la formule (4.2.10)
devient maintenant
\[
 (4.3.1)\qquad\operatorname{Tr}g_{U,\ell}
 =\operatorname{Tr}g_{B_{\unsure{o|0}}^*,\ell}+\operatorname{Tr}g_{T_{\unsure{o|0}}^*,\ell}
 +q^n\operatorname{Tr}g^{-1}_{T_{\unsure{o|0}}^*,\ell}\qquad.
\]
Nous sommes ainsi ramenés à prouver le lemme suivant (où on fera $k\mapsto\bar k'$,
$A\mapsto A_{\unsure{o|0}}^{*\unsure{o|0}}\otimes_{k'}\bar k$) :

\uline{Lemme} 4.3.2. \uline{Soient} $k$ \uline{la clôture algébrique d'un corps fini}, $A$ \uline{un schéma
en groupes commutatif de type fini sur} $k$, $g$ \uline{un automorphisme de ce
schéma en groupes}, \uline{compatible avec l'automorphisme} $f$ \uline{de} $\operatorname{Spec}(k)$ \uline{déduit
par} transport de structure \uline{d'un automorphisme de frobenius itéré}
$x\mapsto x^q$ (\uline{où} $q=p^r$ \uline{est une puissance de la caractéristique}, \uline{avec} $r\geq0$),
\margin{360}
\margin{360}
\margin{172}
